\chapter{Advanced Features}

\section{Cross-referencing}
\LaTeX{} provides excellent cross-referencing capabilities. You can reference:

\begin{itemize}
    \item Sections: See Section~\ref{sec:figures}
    \item Equations: As shown in Equation~\ref{eq:integral}
    \item Figures: Refer to Figure~\ref{fig:example}
    \item Tables: See Table~\ref{tab:comparison} in Chapter 1
\end{itemize}

\section{Figures and Graphics}\label{sec:figures}

\begin{figure}[H]
\centering
% In a real project, you would include an actual image:
% \includegraphics[width=0.8\textwidth]{figures/latex-logo.pdf}
\fbox{\begin{minipage}{0.8\textwidth}
    \centering
    \vspace{2cm}
    \textbf{Placeholder for LaTeX logo}
    \vspace{2cm}
\end{minipage}}
\caption{An example figure showing the LaTeX logo}
\label{fig:example}
\end{figure}

\section{Mathematics}
Complex equations are easy in \LaTeX:

\begin{equation}
\int_{-\infty}^{\infty} e^{-x^2} dx = \sqrt{\pi}
\label{eq:integral}
\end{equation}

Matrix operations:
\[
\begin{bmatrix}
a & b & c \\
d & e & f \\
g & h & i
\end{bmatrix}
\begin{bmatrix}
x \\
y \\
z
\end{bmatrix}
=
\begin{bmatrix}
ax + by + cz \\
dx + ey + fz \\
gx + hy + iz
\end{bmatrix}
\]

\section{Theorems and Proofs}
\begin{theorem}[Pythagorean Theorem]
In a right triangle, the square of the hypotenuse equals the sum of 
the squares of the other two sides.
\end{theorem}

\begin{proof}
Consider a right triangle with sides $a$, $b$, and hypotenuse $c$\ldots
\end{proof}